% Generated draft: requires mathematical and source review.
\documentclass[11pt]{article}
\usepackage[margin=25mm]{geometry}
\usepackage{fontspec}
\setmainfont{Latin Modern Roman}
\usepackage{amsmath,amssymb,mathtools,unicode-math}
\setmathfont{Latin Modern Math}
\usepackage{enumitem,needspace,xurl,hyperref}
\setlength{\parindent}{0pt}
\setlength{\parskip}{5pt}
\setlength{\emergencystretch}{3em}
\newcommand{\N}{\mathbb{N}}
\newcommand{\Z}{\mathbb{Z}}
\newcommand{\Q}{\mathbb{Q}}
\newcommand{\R}{\mathbb{R}}
\newcommand{\C}{\mathbb{C}}
\newcommand{\F}{\mathbb{F}}
\newcommand{\A}{\mathbb{A}}
\newcommand{\PP}{\mathbb{P}}
\newcommand{\E}{\mathbb{E}}
\newcommand{\eps}{\varepsilon}
\newcommand{\dd}{\,\mathrm{d}}
\begin{document}
\section*{UnsolvedMath ID: 30002596}
\textbf{Source problem code:} catalogue entry
\subsection{Definitions and mathematical statement}
For a knot \(K\), write
\[
a(K)=\operatorname{a\!\text{-}\!deg}(K),\qquad p(K)=\operatorname{PL\!\text{-}\!deg}(K),\qquad f(K)=\operatorname{F\!\text{-}\!deg}(K)
\]
for the algebraic, piecewise-linear, and Fourier degrees referred to in the supplied source. These are provisional names: the record does not define the admissible representations, their degrees, or the minimizations defining these knot invariants. It also does not specify the ambient space or equivalence relation for knots. Thus a complete mathematical formulation requires those conventions from the cited report.

Let \(\mathcal K\) denote the intended class of knot types, once specified, and let \(c(K)\) denote the crossing number, namely the minimum number of crossings among admissible regular diagrams of \(K\). The requested problem is to establish comparison inequalities in both directions between each pair among \(a,p,f\), with their constants, hypotheses, and ranges of validity stated explicitly. The source does not specify whether the desired comparisons must be linear, affine, or of some other form.

Conditional on the three degrees being finite and strictly positive on \(\mathcal K\), choose the following orientations for the three pairwise quotients:
\[
r_1(K)=\frac{a(K)}{p(K)},\qquad r_2(K)=\frac{a(K)}{f(K)},\qquad r_3(K)=\frac{p(K)}{f(K)}.
\]
For each \(i\in\{1,2,3\}\), determine the minimum and maximum of \(r_i(K)\) as \(K\) ranges over \(\mathcal K\), or explain any failure of attainment or boundedness by determining the corresponding infimum and supremum. The orientations above are an explicit presentation convention; the source does not select numerator and denominator. Positivity and finiteness are prerequisites for this conditional formulation, not supplied hypotheses verified here.

The asymptotic request concerns limits
\[
\lim_{n\to\infty}r_i(K_n),\qquad c(K_n)\longrightarrow\infty,\qquad i\in\{1,2,3\}.
\]
The record does not say whether a limit is sought for every such sequence, a common limit independent of the sequence, or limits along particular families, nor whether infinite limits are permitted. These alternatives remain unresolved in this draft.

The supplied dated assessment describes one-sided bounds of the forms \(a(K)\leq C f(K)\) and \(a(K)\leq 2p(K)\), and labels the requested reverse comparisons, quotient extrema, and asymptotic limits unresolved. The value and quantification of \(C\), the degree conventions, and the hypotheses of these bounds are absent; these descriptions are retained as supplied background rather than asserted here as verified theorems.
\subsection{Short English statement}
Compare the algebraic, piecewise-linear, and Fourier degrees of knots in both directions, determine the extrema of their pairwise ratios, and investigate their limits as crossing number tends to infinity. The supplied record labels the problem open, but omits the definitions of the degrees and leaves the intended asymptotic quantifiers unclear. This is a conditional statement-only draft requiring mathematical review.
\subsection{Sources}
\begin{itemize}
\item Supplied local UnsolvedMath registry record, permanent ID \(30002596\), source problem code \texttt{catalogue entry}, title: \emph{Relations Among Algebraic PL and Fourier Knot Degrees}; collection: Oberwolfach Reports Open Problems.
\item Supplied source citation: \emph{Algebraic Structures in Low-Dimensional Topology (2015)}, Oberwolfach Reports, pp.\ 1403--1458, DOI \texttt{10.4171/owr/2014/26}. Supplied source URL: \url{https://doi.org/10.4171/owr/2014/26}. The background describes these reports as documenting workshops and their open-problem discussions.
\item Supplied literature-source label: ``Klaus, algebraic, PL and Fourier degrees of knots'', \url{https://doi.org/10.4171/owr/2014/26}. No proposer is identified in the record; this label alone does not establish proposer attribution.
\item Literature checked at: \texttt{2026-08-21}. The open-status assessment and descriptions of known bounds are supplied dated curation derived from the Oberwolfach Reports record, not independently verified here. The supplied curation explicitly requests independent verification of time-sensitive status and cited claims.
\item The record gives proposed year \(2014\), while its report citation includes \(2015\) and its DOI contains \(2014\). These supplied date fields are retained without reconciliation. Its statement status is \texttt{corrected\_verified}; the supplied verification says that mathematical content was unchanged while notation, typography, and wording were normalized. That registry label does not constitute independent mathematical verification of this draft.
\end{itemize}
\end{document}
